% == == == == == == == == == == == == == == == == == == == == == == == == == ==
    == == == == == == == == == == == ==
    = % Section 15 : Preliminary Implementation(Electrical Focus) % == == == ==
      == == == == == == == == == == == == == == == == == == == == == == == == ==
      == == == == == == == == ==
    =

\section {
  Preliminary Implementation
}
\label {
sec:
  implementation
}
\sectionauthors { BS, CM, JH, MN }

\subsection{Implementation Overview}

The development follows three parallel
workstreams(electrical, embedded, software)
that will converge during Capstone 420 integration
    .A modular design philosophy enables subsystems to be developed and tested
        both simultaneously and independently of each other.

\begin{table}[H]
\centering
\caption {
  Implementation Status by Phase
}
\begin{tabular} {
  @{} lll @{}
}
\toprule
\textbf{Phase} & \textbf{Status} & \textbf { Deliverables }
\\
\midrule
\multicolumn{3} { l }
{
  \textit { Capstone 419(Complete) }
}
\ Research &Done &Component selection(power, BMS, sensors,
                                      actuation) \ Design &Done
    &
    6 PCB schematics + layouts(DRC passed) \ Software &Done &ROS 2 architecture,
    communication protocols \ Firmware &Done &Motor control and sensor
    interface framework \\
\midrule
\multicolumn{3} {
  l
}
{
  \textit { Capstone 420(Planned) }
}
\ Build &Pending &PCB fabrication, assembly,
    firmware deployment \ Test &Pending &Unit, integration,
    SLAM / navigation validation \\
\bottomrule
\end {
  tabular
}
\end { table }

\subsection { Electrical Preliminary Implementation }
\label {
sec:
  electrical_impl
}

\subsubsection { Component Selection and Research }

\begin{table}[H]
\centering
\caption { Hardware Subsystem Overview }
\begin{tabular} {
  @{} lll @{}
}
\toprule
\textbf{Subsystem} & \textbf{Function} & \textbf { Key Components }
\\
\midrule Power Input &USB - C charging, power path control &TPS25751D,
    BQ25798 \ Power Distribution &Voltage conversion(6V, 5V, 3.3V) & TPS40305,
    TPS566238, TLV62569 \ Battery Management &Cell monitoring,
    protection, fuel gauge &BQ76920, BQ78350,
    BQ76200 \ Compute \&Control &SLAM / navigation + real
        - time control &Raspberry Pi 5,
    ESP32 - S3 \ Sensing \&Actuation &Environmental perception + motion &LiDAR,
    IMU, motors, drivers \\
\bottomrule
\end {
  tabular
}
\end { table }

\paragraph{Power Input and Charging}
Power enters the system through USB-C, managed by the TPS25751D Dual Role Power (DRP) controller. This enables the robot to either sink power for battery charging or source power to the Raspberry Pi 5.

\begin{table}[H]
\centering
\caption{
  USB - C Power Modes
}
\begin{tabular} {
  @{} lll @{}
}
\toprule
\textbf{Mode} & \textbf{Power} & \textbf { Use Case }
\\
\midrule Sink &Up to 100W(20V / 5A) & Battery charging from adapter \ Source &
    5V / 3A(15W) & Powering Raspberry Pi 5 \\
\bottomrule
\end {
  tabular
}
\end{table}

The BQ25798 buck-boost charger manages 3S Li-ion charging while providing the SYS output for downstream regulators. Buck-boost topology handles USB PD voltages above/below pack voltage seamlessly.

\begin{table}[H]
\centering
\caption{
  BQ25798 Charger Specifications
}
\begin{tabular} {
  @{} ll @{}
}
\toprule
\textbf{Parameter} & \textbf { Value }
\\
\midrule Input voltage &
    3.6 --24V(USB PD 3.0 compatible) \ Charge current &Up to
    5A(10 mA resolution) \ Battery config
    & 3S Li - ion(12.6V charge voltage) \ Peak efficiency &
    96.5\% \ Interface &I\textsuperscript{2} C \\
\bottomrule
\end {
  tabular
}
\end { table }

\paragraph{Power Distribution}
The SYS output feeds three switching regulators that provide isolated rails for different load types.

\begin{table}[H]
\centering
\caption{
  Power Rail Summary
}
\label {
tab:
  power - rails
}
\begin{tabular} {
  @{} lllll @{}
}
\toprule
\textbf{Rail} & \textbf{Regulator} & \textbf{Current} & \textbf{
    Source} & \textbf {
  Loads
}
\\
\midrule 6V & TPS40305DRCT & 25A max &PACK + (AFE) &
    4 $\times$ DC motors \ 5V & TPS566238PRQFR & 6A & SYS &Encoders,
    USB PD \ 3.3V & TLV62569ADRLR & 2A & SYS &ESP32, IMU, sensors \\
\bottomrule
\end {
  tabular
}
\end { table }

\begin{figure}[H]
\centering
\begin{
    tikzpicture}[node distance = 0.6cm and 1cm,
                 block /.style = {rectangle, draw, rounded corners,
                                  minimum width = 1.8cm, minimum height = 0.7cm,
                                  align = center, font =\tiny},
                 rail /.style =
                 {
                   rectangle, draw, fill = #1, minimum width = 1.2cm,
                                    minimum height = 0.5cm, align = center,
                                    font =\tiny\bfseries
                 },
                 arrow /.style = {->, thick, >= stealth}] %
    Battery(top row)
    \node[block, fill = orange !20](gauge){BQ78350\\Fuel Gauge};
\node[block, fill = orange !20, right = 1cm of gauge](afe){BQ76920\\AFE};
\node[block, fill = yellow !30, right = 1cm of afe](battery){3S Li - ion\\Pack};

% Input(bottom row)
    \node[block, fill = gray !30, below = 0.8cm of gauge](usbc){USB -
                                                                 C\\(100W PD)};
\node[block, fill = blue !20, right = 1cm of usbc](drp){TPS25751D\\DRP};
\node[block, fill = green !20, right = 1cm of drp](charger){BQ25798\\Charger};

% SYS rail
    \node[rail = cyan !30, right = 1cm of charger](sys){SYS};

% Power rails
    \node[rail = red !30, right = 1.5cm of sys, yshift = 0.8cm](6v){6V};
\node[rail = blue !30, right = 1.5cm of sys](5v){5V};
\node[rail = green !30, right = 1.5cm of sys, yshift = -0.8cm](3v3){3.3V};

% Loads
    \node[block, fill = red !15, right = 0.8cm of 6v](motors){
      Motors\\(4 $\times$)};
\node[block, fill = blue !15, right = 0.8cm of 5v](pi5){Pi 5\\Encoders};
\node[block, fill = green !15, right = 0.8cm of 3v3](esp32){ESP32\\Sensors};

% Connections
    \draw[arrow](usbc)--(drp);
\draw[arrow](drp)--(charger);
\draw[arrow](charger)--(sys);
\draw[arrow, <->](charger)--(battery);
\draw[arrow, <->](battery)--(afe);
\draw[arrow](afe)--(gauge);

% Power rails from SYS
    \draw[arrow](sys.east)-- ++(0.3, 0) | -(5v.west);
\draw[arrow](sys.east)-- ++(0.3, 0) | -(3v3.west);

% 6V from AFE / PACK + -route from battery right side to 6V rail
    \draw[arrow](battery.east)-- ++(0.3, 0) | -(6v.west);

% Loads
    \draw[arrow](6v)--(motors);
\draw[arrow](5v)--(pi5);
\draw[arrow](3v3)--(esp32);

\end { tikzpicture }
\caption { Power Distribution Architecture }
\label {
fig:
  power - distribution
}
\end { figure }

\begin{table}[H]
\centering
\caption { Power Regulator Details }
\begin{tabular} {
  @{} llp {
    7cm
  }
  @{}
}
\toprule
\textbf{Rail} & \textbf{Key Feature} & \textbf { Notes }
\\
\midrule 6V(TPS40305) & 1.2 MHz, FSS &Voltage mode with feed - forward;
worst - case 12.8A(4 $\times$3 .2A stall) \ 5V(TPS566238)&D
    - CAP3 &No external compensation;
    20.8 / 10.6 m$\Omega$ HS / LS MOSFETs \ 3.3V(TLV62569) & 35 $\mu$A $I_q$ &
    95\% efficiency;
    $ <
    $2 $\mu$A shutdown \\
\bottomrule
\end {
      tabular
    }
\end {
      table
    }

\textbf{Design Evolution(3.3V):
} Originally used LDO from 5V rail, but ESP32-S3 WiFi peaks (500 mA) caused 0.85W dissipation. Switched to dedicated regulator from SYS output.

\paragraph{Battery System}
The BMS employs three ICs: an Analog Front End (AFE) for cell monitoring/protection, a fuel gauge for state-of-charge estimation, and a gate driver for power path control. The 3S Li-ion pack provides 11.1V nominal (12.6V fully charged).

\begin{table}[H]
\centering
\caption{
  BMS IC Summary
}
\label {
tab:
  bms - ics
}
\begin{tabular} {
  @{} llp {
    6cm
  }
  @{}
}
\toprule
\textbf{IC} & \textbf{Part Number} & \textbf { Function } \\
\midrule
AFE & BQ76920PW & Cell voltage monitoring (14-bit ADC), hardware protection (OV/UV/OCD/SCD), cell balancing \\
Fuel Gauge & BQ78350DBT-R1A & CEDV algorithm for SOC/SOH estimation, SMBus interface to host \\
Gate Driver & BQ76200PW & High-side N-FET driver with charge pump, controls CHG/DSG/PCHG paths \\
\bottomrule
\end{
  tabular
}
\end { table }

\subparagraph{AFE Protection Features(BQ76920)} Hardware -
    level protection operates independently of software :

\begin{table}[H]
\centering
\caption {
  BQ76920 Protection Thresholds
}
\begin{tabular} {
  @{} lll @{}
}
\toprule
\textbf{Protection} & \textbf{Threshold} & \textbf { Method }
\\
\midrule Overvoltage(OV) &
    4.25V / cell &Digital comparator \ Undervoltage(UV) &
    2.8V / cell &Digital comparator \ Overcurrent(OCD) & Configurable &
    2 m$\Omega$ sense resistor \ Short Circuit(SCD) &
    Configurable &Fast analog comparator \ Cell Balancing &
    50 --100 mA &Passive resistive \\
\bottomrule
\end {
  tabular
}
\end { table }

\begin{figure}[H]
\centering
\begin{
    tikzpicture}[node distance = 0.5cm and 1cm,
                 block /.style = {rectangle, draw, rounded corners,
                                  minimum width = 1.5cm, minimum height = 0.5cm,
                                  align = center, font =\tiny},
                 decision /.style = {diamond, draw, minimum width = 1cm,
                                     minimum height = 0.5cm, align = center,
                                     font =\tiny, aspect = 2},
                 alert /.style = {rectangle, draw, fill = red !30,
                                  rounded corners, minimum width = 1.2cm,
                                  minimum height = 0.5cm, align = center,
                                  font =\tiny},
                 arrow /.style = {->, thick, >= stealth}] %
    Monitor
    \node[block, fill = green !20](monitor){Cell\\Monitor};
\node[block, fill = blue !20, right = 1cm of monitor](adc){14 - bit\\ADC};

% Comparators
    \node[decision, fill = yellow !30, right = 1cm of adc](ov){V $ > $ 4.25V
                                                                ? };
\node[decision, fill = yellow !30, below = 0.6cm of ov](uv){V $ < $ 2.8V ? };
\node[decision, fill = yellow !30, below = 0.6cm of uv](ocd){I $ > $ OCD ? };
\node[decision, fill = yellow !30, below = 0.6cm of ocd](scd){SCD ? };

% Alerts
    \node[alert, right = 1cm of ov](ov_alert){OV\\Fault};
\node[alert, right = 1cm of uv](uv_alert){UV\\Fault};
\node[alert, right = 1cm of ocd](ocd_alert){OCD\\Fault};
\node[alert, right = 1cm of scd](scd_alert){SCD\\Fault};

% FET control
    \node[block, fill = orange !20, right = 0.8cm of ov_alert,
           yshift = -1.2cm](gate){Gate\\Driver};
\node[block, fill = gray !30, right = 0.8cm of gate](fets){CHG / DSG\\FETs};

% Connections
    \draw[arrow](monitor)--(adc);
\draw[arrow](adc)--(ov);
\draw[arrow](ov)-- node[above, font =\tiny]{Yes}(ov_alert);
\draw[arrow](ov)-- node[right, font =\tiny]{No}(uv);
\draw[arrow](uv)-- node[above, font =\tiny]{Yes}(uv_alert);
\draw[arrow](uv)-- node[right, font =\tiny]{No}(ocd);
\draw[arrow](ocd)-- node[above, font =\tiny]{Yes}(ocd_alert);
\draw[arrow](ocd)-- node[right, font =\tiny]{No}(scd);
\draw[arrow](scd)-- node[above, font =\tiny]{Yes}(scd_alert);

% Alert to gate driver - clean connection
    \coordinate(bus) at([xshift = 0.4cm] ov_alert.east);
\draw(ov_alert.east)--(bus | -ov_alert.east);
\draw(uv_alert.east)--(bus | -uv_alert.east);
\draw(ocd_alert.east)--(bus | -ocd_alert.east);
\draw(scd_alert.east)--(bus | -scd_alert.east);
\draw(bus | -ov_alert.east)--(bus | -scd_alert.east);
\draw[arrow](bus | -gate.west)--(gate.west);

\draw[arrow](gate)--(fets);

\end { tikzpicture }
\caption { BMS Protection Logic Flow }
\label {
fig:
  bms - protection
}
\end { figure }

\subparagraph { Power FET Configuration(BQ76200) }
\begin{table}[H]
\centering
\caption { Power Stage MOSFETs }
\begin{tabular} {
  @{} lll @{}
}
\toprule
\textbf{Function} & \textbf{Part} & \textbf { Key Spec }
\\
\midrule Discharge(DSG) & CSD19531Q5AT & 1.6 m$\Omega$ $R_{DS(on)} $,
    100A \ Charge(CHG) & N - channel &Series with DSG path \ Precharge(PCHG) &
        FDD4141 &P - channel,
    current - limited soft - start \\
\bottomrule
\end {
  tabular
}
\end{table}

N-channel high-side FETs chosen for lower $R_{DS(on)}$ and maintained ground continuity. Integrated charge pump generates 9--14V above PACK+ for gate drive.


\paragraph{Compute and Control Architecture}
The STAR robot employs a distributed compute architecture that separates high-level processing from real-time control. This section describes the platform selection process, the rationale for choosing specific compute modules, and how they integrate to form a cohesive system.

\subparagraph{Platform Selection and Evolution}
The compute platform underwent evolution during Research phase evaluation against project requirements.

\begin{table}[H]
\centering
\caption{
  Compute Platform Evaluation
}
\begin{tabular} {
  @{} llll @{}
}
\toprule
\textbf{Platform} & \textbf{Pros} & \textbf{Cons} & \textbf {
  Decision
}
\\
\midrule PYNQ - Z2(Zynq) & FPGA acceleration &
    32 - bit ARM(ROS 2 needs 64 - bit) & Rejected \ SK - TDA4VM(TI) &
    AI accelerators &
    100W TDP(battery impractical) & Rejected \ Raspberry Pi 5 & 64 - bit,
    5 --10W, \$95 & --& \textbf {
  Selected
}
\\
\bottomrule
\end { tabular }
\end { table }

\subparagraph{Raspberry Pi 5 : Single Board Computer} The Raspberry Pi 5 runs a
    custom Buildroot Linux image and handles complex software stacks,
    while the ESP32 microcontroller executes real - time motor control firmware.

\begin{table}[H]
\centering
\caption {
  Raspberry Pi 5 Specifications
}
\begin{tabular} {
  @{} ll | ll @{}
}
\toprule
\textbf{Spec} & \textbf{Value} & \textbf{Spec} & \textbf { Value }
\\
\midrule Processor &BCM2712(4 $\times$ A76 @2.4 GHz) & Memory &
    8 GB LPDDR4X \ Storage &MicroSD(Buildroot) & WiFi &
    802.11ac dual - band \ USB & 2 $\times$ 3.0,
    2 $\times$ 2.0 & GPIO & 40 - pin(SPI, I\textsuperscript{2} C) \ Power &
        5V / 5A(USB - C from DRP) & Ethernet &Gigabit \\
\bottomrule
\end {
  tabular
}
\end { table }

\textbf{Responsibilities : } SLAM processing, Nav2 navigation,
    sensor fusion(LiDAR / IMU / encoders), web UI hosting,
    velocity commands to ESP32 via SPI.

\subparagraph{
      ESP32 - S3 - WROOM - 1 - N16 : Real - Time Motor Controller
    } Handles time
        - critical operations(motor PWM, sensor polling)
              that cannot tolerate Linux latency variability.

\begin{table}[H]
\centering
\caption {
  ESP32 - S3 Specifications
}
\begin{tabular} {
  @{} ll | ll @{}
}
\toprule
\textbf{Spec} & \textbf{Value} & \textbf{Spec} & \textbf { Value }
\\
\midrule Processor & 2 $\times$ LX7 @240 MHz &SRAM & 512 KB \ Flash &
    16 MB &FPU &Hardware \ USB &Native OTG &MCPWM &Center - aligned PWM \\
\bottomrule
\end {
  tabular
}
\end { table }

\textbf{Why N16 (No PSRAM):} PSRAM (Pseudo Static RAM) variants dedicate GPIO 35--37 to the memory interface. The N16 frees these pins for the second SPI bus to motor drivers. 512 KB internal SRAM is sufficient.

\textbf{Communication:} Pi 5 (controller) $\leftrightarrow$ ESP32 (peripheral) over SPI for velocity commands and sensor data. ESP32 autonomously maintains PWM generation between updates.

\subparagraph{
  Sensor Array
}

\begin{table}[H]
\centering
\caption { Sensor Specifications }
\label {
tab:
  sensors
}
\begin{tabular} {
  @{} llllll @{}
}
\toprule
\textbf{Sensor} & \textbf{Type} & \textbf{Range} & \textbf{FOV} & \textbf{
    Rate} & \textbf {
  Interface
}
\\
\midrule SICK TiM561 & 2D LiDAR & 10 m & 270\textdegree &
    15 Hz &Ethernet \ RPLiDAR C1 & 2D LiDAR & 12 m & 360\textdegree &
    10 Hz &UART \ BNO055 & 9 - axis IMU &
    --&--&100 Hz &I\textsuperscript{2} C \ HC - SR04(7 $\times$) & Ultrasonic &
    2 --400 cm & 15\textdegree & --&GPIO \\
\bottomrule
\end {
  tabular
}
\end { table }

\textbf{LiDAR Selection : } The SICK TiM561(primary)
was donated by SICK Inc., saving \$3,
    500. The RPLiDAR C1(\$65) provides 360\textdegree {}
backup coverage.The TiM561's higher angular resolution (0.33\textdegree{} vs 0.72\textdegree{}) and IP67 rating make it preferred for precision mapping.

\textbf{IMU(BNO055) : } On -
    chip sensor fusion outputs quaternions and Euler angles directly,
    eliminating host -
            side Kalman filter
                complexity.Provides dead reckoning in featureless environments.

\textbf{Ultrasonics(HC - SR04) :
                } Seven sensors around perimeter detect close -
            range obstacles(2 --30 cm) and
        LiDAR - invisible objects(glass).Temperature -
            compensated via DS18B20($v = 331.3 + 0.606T $ m / s)
                .

\paragraph{Actuation and Drive} Four
            -
            wheel differential drive with independently
            controlled DC gear motors and integrated H
            -
            bridge drivers.

\begin{table}[H]
\centering
\caption {
  Motor and Driver Specifications
}
\label {
tab:
  actuation
}
\begin{tabular} {
  @{} ll | ll @{}
}
\toprule
\multicolumn{2} { c | }
{\textbf{Motor(DFRobot FIT0520)}} & \multicolumn{2} { c }
{
  \textbf { Driver(DRV8243SQDGQRQ1) }
}
\\
\midrule Voltage & 6V DC &Max current & 12A / channel \ Speed &
    210 RPM(no - load) & $R_{DS(on)} $ & 98 m$\Omega$ \ Stall torque &
    10 kg$\cdot$cm &PWM freq &Up to 25 kHz \ Stall current &
    3.2A & Interface &SPI + PWM \ Gear ratio & 34 : 1 & Protection &OCP,
    OVP, TSD, UVLO \ Encoder & 341.2 counts / rev &Current sense &IPROPI output \\
\bottomrule
\end {
  tabular
}
\end { table }

\subparagraph { Control Loop Architecture }

\begin{table}[H]
\centering
\caption { Motor Control Data Flow }
\begin{tabular} {
  @{} cll @{}
}
\toprule
\textbf{Step} & \textbf{Component} & \textbf { Function }
\\
\midrule 1 & Pi 5(ROS 2) & Calculate wheel velocities from Nav2 \ 2 &
    SPI bus &Transmit velocity commands \ 3 &
    ESP32(1 kHz) & PID loop : encoder feedback vs.command \ 4 &
    MCPWM &Generate center - aligned PWM \ 5 &
    Encoders &Close velocity loop(341.2 counts / rev) \\
\bottomrule
\end {
  tabular
}
\end { table }

\textbf { Design Rationale: }
144 mm wheels yield 1.1 m / s max speed.AEC -
    Q100 drivers provide fault protection
        .Integrated current sensing enables stall detection.Architecture
            separates timing -
    critical control(ESP32 / FreeRTOS) from navigation(Pi 5 / Linux)
        .

\subsubsection{Schematic Design}

    All schematics were designed in KiCad 9.0.

\paragraph{AFE(Analog Front End)} The BQ76920 AFE schematic provides cell
    voltage monitoring and hardware protection.

\begin{figure}[H]
\centering
\includegraphics[width = 0.95\textwidth, page = 1] {
  figures / schematics / STAR_AFE.pdf
}
\caption { AFE Schematic(BQ76920) }
\label {
fig:
  sch - afe
}
\end { figure }

\paragraph{Buck - Boost Charger} The BQ25798 charger schematic handles USB -
    PD input and battery charging.

\begin{figure}[H]
\centering
\includegraphics[width = 0.95\textwidth, page = 1] {
  figures / schematics / STAR_BUCK_BOOST_CHARGER.pdf
}
\caption { Buck - Boost Charger Schematic(BQ25798) }
\label {
fig:
  sch - charger
}
\end { figure }

\paragraph{MCU}
The ESP32-S3 microcontroller schematic for real-time motor control.

\begin{figure}[H]
\centering
\includegraphics[width=0.95\textwidth,page=1]{
  figures / schematics / STAR_MCU.pdf
}
\caption { MCU Schematic(ESP32 - S3) }
\label {
fig:
  sch - mcu
}
\end { figure }

\paragraph{Motor Driver}
The DRV8243 H-bridge driver schematic for DC motor control.

\begin{figure}[H]
\centering
\includegraphics[width=0.95\textwidth,page=1]{
  figures / schematics / STAR_MOTOR_DRIVER.pdf
}
\caption { Motor Driver Schematic(DRV8243) }
\label {
fig:
  sch - motor
}
\end { figure }

\paragraph{USB - C} The TPS25751D USB - C PD controller schematic.

\begin{figure}[H]
\centering
\includegraphics[width = 0.95\textwidth, page = 1] {
  figures / schematics / STAR_USB - C.pdf
}
\caption { USB - C PD Controller Schematic(TPS25751D) }
\label {
fig:
  sch - usbc
}
\end { figure }

\paragraph{Fuel Gauge}
The BQ78350 fuel gauge schematic for battery state estimation.

\begin{figure}[H]
\centering
\includegraphics[width=0.95\textwidth,page=1]{
  figures / schematics / STAR_FUEL_GAUGE.pdf
}
\caption { Fuel Gauge Schematic(BQ78350) }
\label {
fig:
  sch - fuel
}
\end { figure }

\paragraph{Gate Driver}
The BQ76200 gate driver schematic for power FET control.

\begin{figure}[H]
\centering
\includegraphics[width=0.95\textwidth,page=1]{
  figures / schematics / STAR_GATE_DRIVER.pdf
}
\caption { Gate Driver Schematic(BQ76200) }
\label {
fig:
  sch - gate
}
\end { figure }

\paragraph{3.3V Buck Regulator}
The TLV62569 buck regulator schematic for 3.3V rail.

\begin{figure}[H]
\centering
\includegraphics[width=0.95\textwidth,page=1]{
  figures / schematics / STAR_P3V3_BUCK.pdf
}
\caption { 3.3V Buck Regulator Schematic(TLV62569) }
\label {
fig:
  sch - p3v3
}
\end { figure }

\paragraph{5V Buck Regulator}
The TPS566238 buck regulator schematic for 5V rail.

\begin{figure}[H]
\centering
\includegraphics[width=0.95\textwidth,page=1]{
  figures / schematics / STAR_P5V0_BUCK.pdf
}
\caption { 5V Buck Regulator Schematic(TPS566238) }
\label {
fig:
  sch - p5v0
}
\end { figure }

\paragraph{6V Buck Regulator}
The TPS40305 buck regulator schematic for 6V motor rail.

\begin{figure}[H]
\centering
\includegraphics[width=0.95\textwidth,page=1]{
  figures / schematics / STAR_P6V0_BUCK.pdf
}
\caption { 6V Buck Regulator Schematic(TPS40305) }
\label {
fig:
  sch - p6v0
}
\end { figure }

\paragraph{LiDAR Breakout}
The LiDAR breakout board schematic for sensor interfacing.

\begin{figure}[H]
\centering
\includegraphics[width=0.95\textwidth,page=1]{
  figures / schematics / STAR_LIDAR_BREAKOUT.pdf
}
\caption { LiDAR Breakout Board Schematic }
\label {
fig:
  sch - lidar
}
\end { figure }

\subsubsection{Printed Circuit Board(PCB) Design}

    All PCBs were designed in KiCad 9.0. Most boards use 2 -
    layer stackups; the AFE and Motor Driver use 4-layer stackups for improved signal integrity, current handling, and thermal management. DRC passed on all boards.

\paragraph{AFE PCB (4-Layer)}
The AFE PCB implements the BQ76920 Analog Front End, which provides critical cell voltage monitoring using a 14-bit ADC and hardware protection features such as Overvoltage (OV), Undervoltage (UV), Overcurrent (OCD), and Short Circuit (SCD) detection. This hardware-level protection operates independently of software, ensuring robust battery safety. The 4-layer stackup provides improved signal integrity for the mixed-signal design and better ground plane coverage for noise-sensitive analog measurements.

\begin{figure}[H]
\centering
\includegraphics[width=0.45\textwidth]{
  figures / pcb / STAR_AFE_layer1_top.pdf
}
\hfill
\includegraphics[width = 0.45\textwidth] {
  figures / pcb / STAR_AFE_layer2_in1.pdf
}
\caption { AFE PCB : Layer 1(Top Cu) and Layer 2(In1.Cu) }
\label {
fig:
  pcb - afe - top
}
\end { figure }

\begin{figure}[H]
\centering
\includegraphics[width = 0.45\textwidth] {
  figures / pcb / STAR_AFE_layer3_in2.pdf
}
\hfill
\includegraphics[width = 0.45\textwidth] {
  figures / pcb / STAR_AFE_layer4_bottom.pdf
}
\caption { AFE PCB : Layer 3(In2.Cu) and Layer 4(Bottom Cu) }
\label {
fig:
  pcb - afe - bottom
}
\end { figure }

\paragraph{Buck-Boost Charger PCB}
The Buck-Boost Charger PCB integrates the BQ25798 buck-boost charger, crucial for managing 3S Li-ion battery charging and providing the SYS output for downstream regulators. Its buck-boost topology efficiently handles USB PD voltages, adapting seamlessly whether the input voltage is above or below the battery pack voltage. The PCB is designed to optimize power flow and thermal performance for reliable charging.
\begin{figure}[H]
\centering
\includegraphics[width=0.45\textwidth]{
  figures / pcb / STAR_BUCK_BOOST_CHARGER_top.pdf
}
\hfill
\includegraphics[width = 0.45\textwidth] {
  figures / pcb / STAR_BUCK_BOOST_CHARGER_bottom.pdf
}
\caption { Buck - Boost Charger PCB Layout(Top and Bottom) }
\label {
fig:
  pcb - charger
}
\end { figure }

\paragraph{MCU PCB}
The MCU PCB is built around the ESP32-S3 microcontroller, specifically chosen for its role in real-time motor control. This board handles time-critical operations such as motor PWM generation and sensor polling, tasks that require precise timing and cannot tolerate the latency variability of a Linux-based system. Its design prioritizes efficient signal routing and robust power delivery for the microcontroller.

\paragraph{Motor Driver PCB (4-Layer)}
The Motor Driver PCB features the DRV8243 H-bridge drivers, enabling independent control of the four DC gear motors for the robot's differential drive system. This board utilizes a 4-layer stackup, a critical design choice for improved current handling capabilities and enhanced thermal management, ensuring stable and reliable motor operation under varying load conditions.
The motor driver uses a 4-layer stackup for improved current handling and thermal management.

\begin{figure}[H]
\centering
\includegraphics[width=0.45\textwidth]{
  figures / pcb / STAR_MOTOR_DRIVER_layer1_top.pdf
}
\hfill
\includegraphics[width = 0.45\textwidth] {
  figures / pcb / STAR_MOTOR_DRIVER_layer2_in1.pdf
}
\caption { Motor Driver PCB : Layer 1(Top Cu) and Layer 2(In1.Cu) }
\label {
fig:
  pcb - motor - top
}
\end { figure }

\begin{figure}[H]
\centering
\includegraphics[width = 0.45\textwidth] {
  figures / pcb / STAR_MOTOR_DRIVER_layer3_in2.pdf
}
\hfill
\includegraphics[width = 0.45\textwidth] {
  figures / pcb / STAR_MOTOR_DRIVER_layer4_bottom.pdf
}
\caption { Motor Driver PCB : Layer 3(In2.Cu) and Layer 4(Bottom Cu) }
\label {
fig:
  pcb - motor - bottom
}
\end { figure }

\paragraph{USB-C PCB}
The USB-C PCB incorporates the TPS25751D Dual Role Power (DRP) controller, serving as the primary interface for power input to the system. This controller enables versatile power management, allowing the robot to either sink power for efficient battery charging or source power to the Raspberry Pi 5. The board is designed to handle up to 100W (20V/5A) for charging and provide 5V/3A (15W) when sourcing power.
\begin{figure}[H]
\centering
\includegraphics[width=0.45\textwidth]{
  figures / pcb / STAR_USB - C_top.pdf
}
\hfill
\includegraphics[width = 0.45\textwidth] {
  figures / pcb / STAR_USB - C_bottom.pdf
}
\caption { USB - C PCB Layout(Top and Bottom) }
\label {
fig:
  pcb - usbc
}
\end { figure }

\paragraph{
    Fuel Gauge PCB} The Fuel Gauge PCB is dedicated to accurate battery state -
        of - charge(SOC) and
    state - of - health(SOH) estimation,
    utilizing the BQ78350 fuel gauge IC with its CEDV algorithm.This component
    provides vital battery telemetry via an SMBus interface to the host
    controller,
    offering real
        - time insights into battery performance and remaining capacity
              .The PCB design focuses on precise measurement and reliable
          communication.
\begin{figure}[H]
\centering
\includegraphics[width = 0.45\textwidth] {
  figures / pcb / STAR_FUEL_GAUGE_top.pdf
}
\hfill
\includegraphics[width = 0.45\textwidth] {
  figures / pcb / STAR_FUEL_GAUGE_bottom.pdf
}
\caption { Fuel Gauge PCB Layout(Top and Bottom) }
\label {
fig:
  pcb - fuel
}
\end { figure }

\paragraph{Gate Driver PCB}
The Gate Driver PCB integrates the BQ76200, a specialized high-side N-FET driver with an integrated charge pump. This component is essential for robust power path control, managing the Charge (CHG), Discharge (DSG), and Precharge (PCHG) FETs. The PCB is engineered to provide efficient and reliable switching of these power FETs, crucial for battery protection and power management.
\begin{figure}[H]
\centering
\includegraphics[width=0.45\textwidth]{
  figures / pcb / STAR_GATE_DRIVER_top.pdf
}
\hfill
\includegraphics[width = 0.45\textwidth] {
  figures / pcb / STAR_GATE_DRIVER_bottom.pdf
}
\caption { Gate Driver PCB Layout(Top and Bottom) }
\label {
fig:
  pcb - gate
}
\end { figure }

\paragraph { Power Regulator PCBs }

\textbf{Note:} The 3.3V, 5V, 6V buck regulators and buck-boost charger are currently being redesigned into a single unified power board to reduce complexity and improve integration.

\subparagraph{3.3V Buck PCB}
The 3.3V Buck Regulator PCB utilizes the TLV62569ADRLR buck converter to efficiently generate a stable 3.3V power rail from the SYS output. This rail supplies critical low-power components such as the ESP32 microcontroller, IMU, and various sensors. The PCB is designed for high efficiency and low noise, providing clean power for sensitive digital and analog circuitry.
\begin{figure}[H]
\centering
\includegraphics[width=0.45\textwidth]{
  figures / pcb / STAR_P3V3_BUCK_top.pdf
}
\hfill
\includegraphics[width = 0.45\textwidth] {
  figures / pcb / STAR_P3V3_BUCK_bottom.pdf
}
\caption { 3.3V Buck Regulator PCB Layout(Top and Bottom) }
\label {
fig:
  pcb - p3v3
}
\end { figure }

\subparagraph{5V Buck PCB}
The 5V Buck Regulator PCB integrates the TPS566238PRQFR buck converter, responsible for producing a stable 5V power rail from the SYS output. This 5V rail is vital for powering components like the Raspberry Pi 5, encoders, and other USB PD related circuitry. The design prioritizes high current capability and stable voltage regulation for these essential system elements.
\begin{figure}[H]
\centering
\includegraphics[width=0.45\textwidth]{
  figures / pcb / STAR_P5V0_BUCK_top.pdf
}
\hfill
\includegraphics[width = 0.45\textwidth] {
  figures / pcb / STAR_P5V0_BUCK_bottom.pdf
}
\caption { 5V Buck Regulator PCB Layout(Top and Bottom) }
\label {
fig:
  pcb - p5v0
}
\end { figure }

\subparagraph{6V Buck PCB}
The 6V Buck Regulator PCB employs the TPS40305DRCT buck converter to supply a robust 6V power rail. This rail is primarily dedicated to powering the four DC motors, drawing up to 25A maximum from the PACK+ (AFE) output. The PCB design emphasizes high current delivery, thermal dissipation, and efficient power conversion to meet the demanding requirements of motor actuation.


\paragraph{LiDAR Breakout PCB}
The LiDAR Breakout PCB serves as an essential interface board for integrating the 2D LiDAR sensors into the system, including both the SICK TiM561 (primary) and RPLiDAR C1 (backup). This PCB facilitates the necessary electrical connections and signal conditioning for these sensors, which provide critical data for SLAM and environmental perception. Its design ensures reliable communication and power delivery to the LiDAR units.
\begin{figure}[H]
\centering
\includegraphics[width=0.45\textwidth]{
  figures / pcb / STAR_LIDAR_BREAKOUT_top.pdf
}
\hfill
\includegraphics[width = 0.45\textwidth] {
  figures / pcb / STAR_LIDAR_BREAKOUT_bottom.pdf
}
\caption { LiDAR Breakout PCB Layout(Top and Bottom) }
\label {
fig:
  pcb - lidar
}
\end { figure }

\subsubsection{3D PCB Renders}

    High -
    quality 3D renders were generated from KiCad to visualize
        component placement and board assembly
            .Boards with incomplete PCB
            layouts(MCU, 6V Buck) are excluded pending design updates.

\paragraph {
  Power Management Boards
}

\begin{figure}[H]
\centering
\includegraphics[width = 0.45\textwidth] {
  figures / 3d_renders / STAR_AFE_3d_top.png
}
\hfill
\includegraphics[width = 0.45\textwidth] {
  figures / 3d_renders / STAR_AFE_3d_bottom.png
}
\caption { AFE Board 3D Render(Top and Bottom) }
\label {
fig:
  3d - afe
}
\end { figure }

\begin{figure}[H]
\centering
\includegraphics[width = 0.45\textwidth] {
  figures / 3d_renders / STAR_BUCK_BOOST_CHARGER_3d_top.png
}
\hfill
\includegraphics[width = 0.45\textwidth] {
  figures / 3d_renders / STAR_BUCK_BOOST_CHARGER_3d_bottom.png
}
\caption { Buck - Boost Charger Board 3D Render(Top and Bottom) }
\label {
fig:
  3d - charger
}
\end { figure }

\begin{figure}[H]
\centering
\includegraphics[width = 0.45\textwidth] {
  figures / 3d_renders / STAR_FUEL_GAUGE_3d_top.png
}
\hfill
\includegraphics[width = 0.45\textwidth] {
  figures / 3d_renders / STAR_FUEL_GAUGE_3d_bottom.png
}
\caption { Fuel Gauge Board 3D Render(Top and Bottom) }
\label {
fig:
  3d - fuel
}
\end { figure }

\begin{figure}[H]
\centering
\includegraphics[width = 0.45\textwidth] {
  figures / 3d_renders / STAR_GATE_DRIVER_3d_top.png
}
\hfill
\includegraphics[width = 0.45\textwidth] {
  figures / 3d_renders / STAR_GATE_DRIVER_3d_bottom.png
}
\caption { Gate Driver Board 3D Render(Top and Bottom) }
\label {
fig:
  3d - gate
}
\end { figure }

\paragraph { Control and Interface Boards }

\begin{figure}[H]
\centering
\includegraphics[width = 0.45\textwidth] {
  figures / 3d_renders / STAR_MOTOR_DRIVER_3d_top.png
}
\hfill
\includegraphics[width = 0.45\textwidth] {
  figures / 3d_renders / STAR_MOTOR_DRIVER_3d_bottom.png
}
\caption { Motor Driver Board 3D Render(Top and Bottom) }
\label {
fig:
  3d - motor
}
\end { figure }

\begin{figure}[H]
\centering
\includegraphics[width = 0.45\textwidth] {
  figures / 3d_renders / STAR_USB - C_3d_top.png
}
\hfill
\includegraphics[width = 0.45\textwidth] {
  figures / 3d_renders / STAR_USB - C_3d_bottom.png
}
\caption { USB - C Board 3D Render(Top and Bottom) }
\label {
fig:
  3d - usbc
}
\end { figure }

\begin{figure}[H]
\centering
\includegraphics[width = 0.45\textwidth] {
  figures / 3d_renders / STAR_LIDAR_BREAKOUT_3d_top.png
}
\hfill
\includegraphics[width = 0.45\textwidth] {
  figures / 3d_renders / STAR_LIDAR_BREAKOUT_3d_bottom.png
}
\caption { LiDAR Breakout Board 3D Render(Top and Bottom) }
\label {
fig:
  3d - lidar
}
\end { figure }

\paragraph { Power Regulator Boards }

\textbf{Note : } These boards are being consolidated into a unified power board.

\begin{figure}[H]
\centering
\includegraphics[width = 0.45\textwidth] {
  figures / 3d_renders / STAR_P3V3_BUCK_3d_top.png
}
\hfill
\includegraphics[width = 0.45\textwidth] {
  figures / 3d_renders / STAR_P3V3_BUCK_3d_bottom.png
}
\caption { 3.3V Buck Regulator Board 3D Render(Top and Bottom) }
\label {
fig:
  3d - p3v3
}
\end { figure }

\begin{figure}[H]
\centering
\includegraphics[width = 0.45\textwidth] {
  figures / 3d_renders / STAR_P5V0_BUCK_3d_top.png
}
\hfill
\includegraphics[width = 0.45\textwidth] {
  figures / 3d_renders / STAR_P5V0_BUCK_3d_bottom.png
}
\caption { 5V Buck Regulator Board 3D Render(Top and Bottom) }
\label {
fig:
  3d - p5v0
}
\end { figure }

\subsubsection{Bill of Materials}

Complete BOMs were exported from KiCad 9.0 for each PCB with the following fields:
\begin{
  itemize
}
\item Reference designator, Value, Footprint,
    Quantity
    \item Manufacturer name and part number(MPN)
    \item Digi - Key part number and unit price
    \item Mouser part number and unit price
    \item Component description
\end{itemize}

                  Full BOMs are provided in CSV format in the project
                  repository(\texttt{figures / bom\_full / })
                      .

\paragraph {
  PCB Layer Stackup Summary
}
\begin{table}[H]
\centering
\caption { PCB Layer Configuration }
\label {
tab:
  layer - stackup
}
\begin{tabular} {
  @{} llll @{}
}
\toprule
\textbf{Board} & \textbf{Layers} & \textbf{Stackup} & \textbf {
  Rationale
}
\\
\midrule AFE & 4 & F.Cu / In1.Cu / In2.Cu / B.Cu &Mixed - signal integrity,
    ground plane coverage \ Motor Driver & 4 &
        F.Cu / In1.Cu / In2.Cu / B.Cu &High current(12A),
    thermal management \\
\midrule Buck - Boost Charger & 2 & F.Cu / B.Cu &Power stage,
    adequate copper area \ Fuel Gauge & 2 & F.Cu / B.Cu &Digital + analog,
    low power \ Gate Driver & 2 & F.Cu / B.Cu &Simple gate drive circuit \ MCU &
        2 & F.Cu / B.Cu &Digital signals,
    connectors \ USB - C & 2 & F.Cu / B.Cu &USB PD controller,
    standard routing \ 3.3V Buck &
        2 & F.Cu / B.Cu &Low - power regulator \ 5V Buck &
        2 & F.Cu / B.Cu &Medium - power regulator \ 6V Buck &
        2 & F.Cu / B.Cu &Motor power regulator \ LiDAR Breakout &
        2 & F.Cu / B.Cu &Simple interface board \\
\bottomrule
\end {
  tabular
}
\end { table }

\paragraph{Breadboard Testing Limitations}

Breadboards cannot be used for prototyping or testing the power electronics in this project. The 6V motor rail must deliver up to 25A peak current (four motors at 3.2A stall each), while standard breadboard contact clips are rated for approximately 1A. Exceeding this limit causes resistive heating at the contacts, leading to unreliable connections and potential fire hazards. This current limitation necessitated moving directly to custom PCB fabrication for the power distribution, motor driver, and battery management subsystems. Phase~1 hardware verification in Section~\ref{sec:test_plan} (Table~\ref{tab:phase1_summary}) therefore focuses on PCB-level tests (TC-PCB-001 through TC-PCB-006) to verify solder quality and trace continuity under load conditions that breadboards cannot safely support.

\paragraph{
  Component Count Summary
}
\begin{table}[H]
\centering
\caption { BOM Summary by Board }
\label {
tab:
  bom - summary
}
\begin{tabular} {
  @{} lrrll @{}
}
\toprule
\textbf{Board} & \textbf{Layers} & \textbf{Components} & \textbf{
    Primary IC} & \textbf {
  Package
}
\\
\midrule AFE & 4 & 81 & BQ76920PW &TSSOP - 20 \ Buck - Boost Charger & 2 & 41 &
    BQ25798RQMR &QFN - 32 \ Fuel Gauge & 2 & 46 &
    BQ78350DBT - R1A &TSSOP - 30 \ Gate Driver & 2 & 12 &
    BQ76200PW &TSSOP - 14 \ MCU & 2 & 42 &
    ESP32 - S3 - WROOM - 1 - N16 &Module \ Motor Driver & 4 & 42 &
    DRV8243SQDGQRQ1(4 $\times$) & WQFN - 29 \ USB - C & 2 &
    49 & TPS25751DREFR &QFN - 64 \ 3.3V Buck & 2 &
    15 & TLV62569ADRLR &SOT - 5X3 \ 5V Buck & 2 &
    23 & TPS566238PRQFR &QFN - 18 \ 6V Buck & 2 &
    34 & TPS40305DRCT &SON - 10 \ LiDAR Breakout & 2 & 5 & --&-- \\
\midrule
\textbf{Total} & --& \textbf{390} & --&-- \\
\bottomrule
\end {
  tabular
}
\end { table }

\paragraph { Motor Driver BOM(42 components) }
\begin{table}[H]
\centering
\caption { Motor Driver Board BOM with Pricing }
\label {
tab:
  bom - motor - full
}
\footnotesize
\begin{tabular} {
  @{} llrlrr @{}
}
\toprule
\textbf{Ref} & \textbf{Value} & \textbf{Qty} & \textbf{MPN} & \textbf{
    DigiKey} & \textbf {
  Mouser
}
\\
\midrule U1-- U4 &DRV8243S & 4 &
    DRV8243SQDGQRQ1 & \$4 .72 & \$6 .37 \ C1-- C4 & 56 $\mu$F & 4 &
    865080343007 & \$0 .26 & \$0 .27 \ C5-- C8 & 220nF & 4 &
    CL10B224KB8NNNC & \$0 .11 & \$0 .11 \ C9-- C12 & 100nF & 4 &
    CC603JRX7R9BB104 & \$0 .10 & \$0 .10 \ R1-- R4 & 10k $\Omega$ & 4 &
    RT0603BRE0710KL & \$0 .14 & \$0 .10 \ R5-- R8 & 68 $\Omega$ & 4 &
    RT0603FRE0768RL & \$0 .10 & \$0 .10 \ R9-- R12 & 3.16k $\Omega$ & 4 &
    ERA - 3AEB3161V & \$0 .10 & \$0 .10 \ D1-- D4 &Red LED & 4 &
    150060RS75000 & \$0 .15 & \$0 .16 \ J3,
    J10-- 12 & B6B - XH - A & 4 & B6B - XH - A(JST) & --&-- \ J7 &Phoenix 2P &
        1 & 1017491 & --&-- \\
\bottomrule
\end {
  tabular
}
\end { table }

\paragraph { USB - C PD Controller BOM(49 components) }
\begin{table}[H]
\centering
\caption { USB - C Board BOM with Pricing }
\label {
tab:
  bom - usbc - full
}
\footnotesize
\begin{tabular} {
  @{} llrlrr @{}
}
\toprule
\textbf{Ref} & \textbf{Value} & \textbf{Qty} & \textbf{MPN} & \textbf{
    DigiKey} & \textbf {
  Mouser
}
\\
\midrule U2 &TPS25751D & 1 & TPS25751DREFR & \$2 .90 & \$2 .91 \ U1 &EEPROM &
    1 & CAT24C512WI - GT3 & \$0 .59 & \$0 .59 \ U3 &TVS & 1 &
    TVS2200DRVR & \$0 .53 & \$0 .83 \ J7 &USB - C & 1 &
    USB4216 - 03 - A & \$0 .62 & \$0 .72 \ C3-- C13 & 10 $\mu$F & 11 &
    GRM21BR61H106KE3L & \$0 .28 & \$0 .28 \ C23 & 22 $\mu$F & 1 &
    T527T226M025ATE100 &N / A & \$1 .13 \ C26-- C31 & 4.7 $\mu$F & 4 &
    GRT21BR61H475KE3K & \$0 .17 & \$0 .17 \ D1,
    D2 &TVS & 2 & 82306050029(Würth) & \$0 .73 & \$0 .73 \ R1-- R4 &
        10k $\Omega$ & 4 & RT0603BRE0710KL & \$0 .14 & \$0 .10 \\
\bottomrule
\end {
  tabular
}
\end { table }

\paragraph { MCU Board BOM(42 components) }
\begin{table}[H]
\centering
\caption { MCU Board BOM with Pricing }
\label {
tab:
  bom - mcu - full
}
\footnotesize
\begin{tabular} {
  @{} llrlrr @{}
}
\toprule
\textbf{Ref} & \textbf{Value} & \textbf{Qty} & \textbf{MPN} & \textbf{
    DigiKey} & \textbf {
  Mouser
}
\\
\midrule U2 &ESP32 - S3 & 1 &
    ESP32 - S3 - WROOM - 1 - N16 & \$3 .50 & \$3 .50 \ U1 &Temp Sensor & 1 &
    DS18B20 + & \$4 .25 & \$4 .25 \ U3 &Decoder & 1 & 74HC139D,
    653 & \$0 .45 & \$0 .45 \ U4 &Mux & 1 &
        SN74HC153DR & \$0 .35 & \$0 .35 \ J1-- J7 &HC - SR04 & 7 &
        4 - pin header & \$0 .10 & \$0 .10 \ J8 &Pi GPIO & 1 &
        2x20 header & \$1 .50 & \$1 .50 \ D1,
    D2 &ESD & 2 & 82306050029 & \$0 .73 & \$0 .73 \ R1-- R8 & 10k $\Omega$ & 8 &
        RT0603BRE0710KL & \$0 .14 & \$0 .10 \\
\bottomrule
\end {
  tabular
}
\end { table }

\paragraph { Power Management BOMs }
\begin{table}[H]
\centering
\caption {
  Power Board BOM Summary with Key Component Pricing
}
\label {
tab:
  bom - power
}
\footnotesize
\begin{tabular} {
  @{} llrlrr @{}
}
\toprule
\textbf{Board} & \textbf{Primary IC} & \textbf{Qty} & \textbf{MPN} & \textbf{
    DigiKey} & \textbf {
  Mouser
}
\\
\midrule Buck - Boost Charger &BQ25798 & 41 &
    BQ25798RQMR & \$5 .50 & \$5 .50 \ Fuel Gauge &BQ78350 & 46 &
    BQ78350DBT - R1A & \$8 .25 & \$8 .25 \ Gate Driver &BQ76200 & 12 &
    BQ76200PW & \$2 .15 & \$2 .15 \ AFE &BQ76920 & 81 &
    BQ7692000PWR & \$3 .75 & \$3 .75 \ 3.3V Buck &TLV62569 & 15 &
    TLV62569ADRLR & \$1 .25 & \$1 .25 \ 5V Buck &TPS566238 & 23 &
    TPS566238PRQFR & \$1 .85 & \$1 .85 \ 6V Buck &TPS40305 & 34 &
    TPS40305DRCT & \$2 .45 & \$2 .45 \\
\bottomrule
\end {
  tabular
}
\end{table}

Complete BOMs with all component fields for each PCB are provided in Appendix~\ref{app:bom}.
